\documentclass[11pt]{amsart}
\usepackage{amssymb}
\usepackage[margin=1.15in]{geometry}
\usepackage[hidelinks]{hyperref}
\hypersetup{pdftitle={Step II over an arbitrary coefficient ring}}

\newtheorem{theorem}{Theorem}[section]
\newtheorem{lemma}[theorem]{Lemma}
\theoremstyle{definition}
\newtheorem{definition}[theorem]{Definition}
\theoremstyle{remark}
\newtheorem{remark}[theorem]{Remark}
\newcommand{\code}[1]{\texttt{#1}}
\newcommand{\mere}[1]{\left\lVert #1\right\rVert}

\title{Step II over an arbitrary coefficient ring}
\date{October 26, 2026}

\begin{document}
\begin{abstract}
Let $A$ be any commutative ring, $F,G,U,V\in A[X,Y]$ with $FU=GV$, and
$J$ the Jacobian of $F,G$. We show $J^eU\in(G)$ for some $e$, without
the annihilator stabilization hypothesis of
\code{jacobian-regularity.pdf} and without passing to a noetherian
coefficient ring by a classical argument. Noetherianity is used once,
along one sequence, in the constructive form of Coquand--Persson, inside
a subring of $B=(A[X,Y]/(G))[1/J]$ that is given by codes.
\end{abstract}
\maketitle

Throughout, $B=(A[X,Y]/(G))[1/J]$, and $x,y,f,u$ are the images of
$X,Y,F,U$ in $B$. As in \code{jacobian-regularity.pdf}, it suffices to
prove $u=0$ in $B$: the kernels of localization and of the quotient by
$G$ then give $J^eU=GW$ in $A[X,Y]$.

\section{The kill argument}

For each $n$, finite Taylor lifting gives $x_n(T),y_n(T)\in B[T]$ with
$x_n(0)=x$, $y_n(0)=y$, $F(x_n,y_n)\equiv f-T$ and $G(x_n,y_n)\equiv0$
modulo $T^{n+1}$. Substituting in $FU=GV$ gives
$(f-T)\,p_n(T)\equiv0\pmod{T^{n+1}}$ with $p_n=U(x_n,y_n)$. Comparing
coefficients, with $w_n$ the coefficient of $T^n$ in $p_n$,
\[
 u=p_n(0)=f^n w_n,\qquad f^{n+1}w_n=0 .
\]

\begin{lemma}[Kill]\label{lem:kill}
Let $u,f,w_0,w_1,\ldots$ be elements of a ring with $u=f^nw_n$ and
$f^{n+1}w_n=0$ for all $n$. If some $w_k$ is a linear combination of
$w_0,\ldots,w_{k-1}$, then $u=0$.
\end{lemma}
\begin{proof}
For $j<k$, $f^kw_j=f^{k-1-j}f^{j+1}w_j=0$, so $f^k$ kills $w_k$, and
$u=f^kw_k=0$. For $k=0$ the empty combination gives $w_0=0=u$.
\end{proof}

So it suffices that the sequence $(w_n)$ has a \emph{good} initial
segment: one whose last element is a combination of the earlier ones.
In a noetherian ring in the sense of Coquand--Persson (``good bars the
empty list'') every sequence has one, provided the elements of the
sequence can be followed along the tree of the bar; this needs a code
for each $w_n$, but only merely, since the conclusion $u=0$ is a
proposition.

\section{The coded subring}

Let $a_1,\ldots,a_m$ be the coefficients of $F,G,U,V$, and let
$A_0\subseteq A$ be the subring they generate. Its image in $B$,
together with $x$, $y$ and $1/J$, generates a subring of $B$ that is a
quotient of $\mathbb Z[t_1,\ldots,t_{m+3}]$; we present it by codes.

\begin{definition}
Let $K$ be the type of codes of $\mathbb Z[t_1,\ldots,t_{m+3}]$
(iterated polynomial codes over the codes of $\mathbb Z$), and
$\varphi:K\to B$ the evaluation at $h(a_1),\ldots,h(a_m),x,y,1/J$. It
preserves the ring operations; no law is asked of $K$. The \emph{coded
subring} is
\[
 B'=\sum_{b:B}\ \mere{\textstyle\sum_{c:K}\varphi(c)=b}.
\]
\end{definition}

The predicate is a proposition, so $B'$ is a set-level $\Sigma$ and the
projection $\iota:B'\to B$ is injective: elements of $B'$ are equal
when their images are. Coded elements are closed under the ring
operations (apply the operation to codes, under the truncation), so
$B'$ is a commutative ring, every law holding because it holds in $B$,
and $\iota$ is a ring map. The elements $h(a_i)$, $x$, $y$, $1/J$ lie
in $B'$, with the codes $t_i$.

\begin{lemma}\label{lem:coded}
$B'$ is \emph{coded noetherian}: the presentation $\varphi':K\to B'$,
$\varphi'(c)=(\varphi(c),\mere{(c,\mathrm{refl})})$, is noetherian, and
every element of $B'$ merely has a code.
\end{lemma}
\begin{proof}
The second part is the definition of $B'$. For the first, $\varphi$ is
noetherian by the Hilbert basis theorem iterated $m+3$ times over the
image of $\mathbb Z$ (\code{sag/noetherian}). A combination in $B'$ with
coefficients in $K$ is one in $B$ with the same coefficients, and
conversely, since $\iota$ is injective and $\iota\circ\varphi'=\varphi$.
So the good lists for $\varphi'$ and $\varphi$ are the same lists of
codes, and the tree of the bar for $\varphi$ is one for $\varphi'$.
\end{proof}

\section{Running the lifting in \texorpdfstring{$B'$}{B'}}

Write $F',G',U',V'$ for $F,G,U,V$ with coefficients mapped to $B'$
(through $A_0\to B'$). In $B'$ we have $G'(x,y)=0$ and $J'(x,y)$
invertible with inverse $1/J$: both equations hold after $\iota$, hence
in $B'$. So Taylor lifting can be run in $B'$, giving
$x'_n,y'_n\in B'[T]$, $p'_n=U'(x'_n,y'_n)$ and
\[
 u'=f'^{\,n}w'_n,\qquad f'^{\,n+1}w'_n=0\qquad\text{in }B',
\]
with $\iota(u')=u$. The argument of the previous section uses the
relation $FU=GV$ only through its consequence
\[
 F'(x'_n,y'_n)\,U'(x'_n,y'_n)=G'(x'_n,y'_n)\,V'(x'_n,y'_n)
 \qquad\text{in }B'[T].
\]
This is not obtained from a relation in $A_0[X,Y]$. Apply $\iota$ to
coefficients: $\iota[T]:B'[T]\to B[T]$ is a ring map, and carries
$H'(x'_n,y'_n)$ to $H(\iota x'_n,\iota y'_n)$ for every polynomial $H$
over $A$. In $B[T]$ the identity holds, as the image of $FU=GV$ under
the $A$-algebra map $A[X,Y]\to B[T]$, $X\mapsto\iota x'_n$,
$Y\mapsto\iota y'_n$. Since $\iota$ is injective, the two sides agree
coefficientwise in $B'$.

\section{Conclusion}

Each $w'_n$ is an element of $B'$, so it merely has a code: this is the
point of running the lifting in $B'$ rather than in $B$, where the
coefficients of the lift are built from $x,y$, $1/J$ and the $a_i$ by
the ring operations but this is not recorded. By Lemma~\ref{lem:coded}
the walk along the tree of the bar, following merely given codes of
$w'_0,w'_1,\ldots$, merely finds a good initial segment, and
Lemma~\ref{lem:kill} gives $u'=0$, a proposition, so the truncation is
eliminated without choice. Then $u=\iota(u')=0$ in $B$.

\begin{theorem}
For every commutative ring $A$ and $F,G,U,V\in A[X,Y]$ with $FU=GV$,
merely $J^eU=GW$ for some $e\geq0$ and $W\in A[X,Y]$.
\end{theorem}

No excluded middle, choice, decidable equality, or noetherian
hypothesis on $A$ is used: the only noetherian rings are the finitely
generated ones over $\mathbb Z$, presented by codes.

\section*{Formalisation}

In \code{unit/sag/jacobian}:
\code{CodedSubring} (the subring $B'$ of a presentation, its ring
structure, Lemma~\ref{lem:coded}, and Lemma~\ref{lem:kill} along a
merely coded sequence, from \code{Sequence.Kill.Walk});
\code{RawMap} (a ring map on coefficient lists);
\code{RegularAt.Curve} (the order-$n$ calculation from the curve
identity alone);
\code{FiniteCoefficients} (the generators, $F',G',U',V'$, the transfer
of $G(x,y)=0$, of $1/J$ and of the curve identity along $\iota$, and the
theorem). In the formal text the coefficient map is the identity of
$B'$: only the images $h(a_i)$ of the generators of $A_0$ are needed,
so $A_0$ is not built as a ring of its own.
\end{document}
